% =========================================================
% Proof: Finite Union of Closed Sets is Closed
% Mode: standard
% =========================================================

\subsection*{Finite Union of Closed Sets is Closed}
\label{prf:finite-union-closed}

\begin{proof}
\Claim Let $(X,d)$ be a metric space and let $F_1, \dots, F_n \subseteq X$ be closed.
Then $\bigcup_{k=1}^n F_k$ is closed.

\Given Closed sets $F_1, \dots, F_n \subseteq X$.

\Goal To show $F := \bigcup_{k=1}^n F_k$ is closed, i.e.\ $X \setminus F$ is open.

\Strategy Use the complement characterization and De Morgan's law.
The complement of a finite union is a finite intersection of open sets,
and finite intersections of open sets are open.

\medskip
\ByAss each $F_k$ is closed, so each $X \setminus F_k$ is open.

\ByThm{De Morgan's law}:
\[
X \setminus F \;=\; X \setminus \bigcup_{k=1}^n F_k \;=\; \bigcap_{k=1}^n (X \setminus F_k).
\]

\ByThm{Proposition~\ref{prop:finite-intersection-open}} a finite intersection of open sets is open.

\Therefore $X \setminus F$ is open, so $F$ is closed. \AsReq
\end{proof}

\begin{remark*}[Why finiteness is essential]
An infinite union of closed sets need not be closed.
For example in $\mathbb{R}$: $\{1/n\}$ is closed for each $n \ge 1$,
but $\bigcup_{n=1}^\infty \{1/n\} = \{1, 1/2, 1/3, \ldots\}$
is not closed: $0$ is a limit point not in the set.
\end{remark*}

\begin{remark*}[Dependencies]
\begin{itemize}
  \item Definition~\ref{def:open-closed-set}: $F$ closed $\iff$ $X \setminus F$ open.
  \item De Morgan's law: $X \setminus \bigcup_k F_k = \bigcap_k (X \setminus F_k)$.
  \item Proposition~\ref{prop:finite-intersection-open}: finite intersections of open sets are open.
\end{itemize}
\end{remark*}
